\subsection{Worked solutions on the ten soils}
The steady flux between two depths depends on the conductivity at the suction in between. This section computes it for each of the ten soils.

\exbox{Ten-soil example 2a: Coarse loam (S01)}{Coarse loam: $K_s=31.2$ cm\,d$^{-1}$ (typical for the class), $\alpha=0.0961$, $n=1.353$. Find $K(h)$ at 100, 300 and 1000 cm suction and the steady upward flux between $h=300$ cm at 0.4 m and $h=100$ cm at 0 m, using Darcy's law $q=K(\Delta h/\Delta z-1)$ with the geometric mean of $K$.}\sol $K_r(100)=9.39e-05$, $K_r(300)=4.16e-06$, $K_r(1000)=1.31e-07$ \ans $K=0.00293$, $0.00013$, $4.08e-06$ cm\,d$^{-1}$ \ans Geometric mean of the two ends: $\bar K=\sqrt{0.00293\times0.00013}=0.000616$ cm\,d$^{-1}$ \ans Gradient: $\Delta h/\Delta z=(300-100)/40=5.0$, so $q=0.000616\times(5.0-1)=0.00246$ cm\,d$^{-1}$ \ans The flux is negligible against the demand of a crop (0.5 cm per day).

\exbox{Ten-soil example 2a: Coarse sand (S02)}{Coarse sand: $K_s=288.0$ cm\,d$^{-1}$ (typical for the class), $\alpha=0.0536$, $n=1.600$. Find $K(h)$ at 100, 300 and 1000 cm suction and the steady upward flux between $h=300$ cm at 0.4 m and $h=100$ cm at 0 m, using Darcy's law $q=K(\Delta h/\Delta z-1)$ with the geometric mean of $K$.}\sol $K_r(100)=0.000356$, $K_r(300)=8.28e-06$, $K_r(1000)=1.24e-07$ \ans $K=0.102$, $0.00239$, $3.58e-05$ cm\,d$^{-1}$ \ans Geometric mean of the two ends: $\bar K=\sqrt{0.102\times0.00239}=0.0156$ cm\,d$^{-1}$ \ans Gradient: $\Delta h/\Delta z=(300-100)/40=5.0$, so $q=0.0156\times(5.0-1)=0.0625$ cm\,d$^{-1}$ \ans The flux is negligible against the demand of a crop (0.5 cm per day).

\exbox{Ten-soil example 2a: Sand (S03)}{Sand: $K_s=504.0$ cm\,d$^{-1}$ (typical for the class), $\alpha=0.1450$, $n=2.680$. Find $K(h)$ at 100, 300 and 1000 cm suction and the steady upward flux between $h=300$ cm at 0.4 m and $h=100$ cm at 0 m, using Darcy's law $q=K(\Delta h/\Delta z-1)$ with the geometric mean of $K$.}\sol $K_r(100)=2.47e-08$, $K_r(300)=2.73e-11$, $K_r(1000)=1.56e-14$ \ans $K=1.25e-05$, $1.37e-08$, $7.88e-12$ cm\,d$^{-1}$ \ans Geometric mean of the two ends: $\bar K=\sqrt{1.25e-05\times1.37e-08}=4.14e-07$ cm\,d$^{-1}$ \ans Gradient: $\Delta h/\Delta z=(300-100)/40=5.0$, so $q=4.14e-07\times(5.0-1)=1.66e-06$ cm\,d$^{-1}$ \ans The flux is negligible against the demand of a crop (0.5 cm per day).

\exbox{Ten-soil example 2a: Loamy sand (S04)}{Loamy sand: $K_s=146.4$ cm\,d$^{-1}$ (typical for the class), $\alpha=0.1240$, $n=2.280$. Find $K(h)$ at 100, 300 and 1000 cm suction and the steady upward flux between $h=300$ cm at 0.4 m and $h=100$ cm at 0 m, using Darcy's law $q=K(\Delta h/\Delta z-1)$ with the geometric mean of $K$.}\sol $K_r(100)=6.46e-07$, $K_r(300)=2.15e-09$, $K_r(1000)=4.1e-12$ \ans $K=9.46e-05$, $3.14e-07$, $6e-10$ cm\,d$^{-1}$ \ans Geometric mean of the two ends: $\bar K=\sqrt{9.46e-05\times3.14e-07}=5.45e-06$ cm\,d$^{-1}$ \ans Gradient: $\Delta h/\Delta z=(300-100)/40=5.0$, so $q=5.45e-06\times(5.0-1)=2.18e-05$ cm\,d$^{-1}$ \ans The flux is negligible against the demand of a crop (0.5 cm per day).

\exbox{Ten-soil example 2a: Sandy loam (S05)}{Sandy loam: $K_s=62.4$ cm\,d$^{-1}$ (typical for the class), $\alpha=0.0750$, $n=1.890$. Find $K(h)$ at 100, 300 and 1000 cm suction and the steady upward flux between $h=300$ cm at 0.4 m and $h=100$ cm at 0 m, using Darcy's law $q=K(\Delta h/\Delta z-1)$ with the geometric mean of $K$.}\sol $K_r(100)=4.29e-05$, $K_r(300)=4.27e-07$, $K_r(1000)=2.65e-09$ \ans $K=0.00268$, $2.67e-05$, $1.65e-07$ cm\,d$^{-1}$ \ans Geometric mean of the two ends: $\bar K=\sqrt{0.00268\times2.67e-05}=0.000267$ cm\,d$^{-1}$ \ans Gradient: $\Delta h/\Delta z=(300-100)/40=5.0$, so $q=0.000267\times(5.0-1)=0.00107$ cm\,d$^{-1}$ \ans The flux is negligible against the demand of a crop (0.5 cm per day).

\exbox{Ten-soil example 2a: Loam (S06)}{Loam: $K_s=31.2$ cm\,d$^{-1}$ (typical for the class), $\alpha=0.0360$, $n=1.560$. Find $K(h)$ at 100, 300 and 1000 cm suction and the steady upward flux between $h=300$ cm at 0.4 m and $h=100$ cm at 0 m, using Darcy's law $q=K(\Delta h/\Delta z-1)$ with the geometric mean of $K$.}\sol $K_r(100)=0.00136$, $K_r(300)=3.8e-05$, $K_r(1000)=6.55e-07$ \ans $K=0.0424$, $0.00119$, $2.04e-05$ cm\,d$^{-1}$ \ans Geometric mean of the two ends: $\bar K=\sqrt{0.0424\times0.00119}=0.00709$ cm\,d$^{-1}$ \ans Gradient: $\Delta h/\Delta z=(300-100)/40=5.0$, so $q=0.00709\times(5.0-1)=0.0284$ cm\,d$^{-1}$ \ans The flux is negligible against the demand of a crop (0.5 cm per day).

\exbox{Ten-soil example 2a: Silt loam (S07)}{Silt loam: $K_s=16.3$ cm\,d$^{-1}$ (typical for the class), $\alpha=0.0200$, $n=1.410$. Find $K(h)$ at 100, 300 and 1000 cm suction and the steady upward flux between $h=300$ cm at 0.4 m and $h=100$ cm at 0 m, using Darcy's law $q=K(\Delta h/\Delta z-1)$ with the geometric mean of $K$.}\sol $K_r(100)=0.00652$, $K_r(300)=0.000335$, $K_r(1000)=9.6e-06$ \ans $K=0.106$, $0.00547$, $0.000157$ cm\,d$^{-1}$ \ans Geometric mean of the two ends: $\bar K=\sqrt{0.106\times0.00547}=0.0241$ cm\,d$^{-1}$ \ans Gradient: $\Delta h/\Delta z=(300-100)/40=5.0$, so $q=0.0241\times(5.0-1)=0.0964$ cm\,d$^{-1}$ \ans The flux is negligible against the demand of a crop (0.5 cm per day).

\exbox{Ten-soil example 2a: Sandy clay loam (S08)}{Sandy clay loam: $K_s=10.3$ cm\,d$^{-1}$ (typical for the class), $\alpha=0.0590$, $n=1.480$. Find $K(h)$ at 100, 300 and 1000 cm suction and the steady upward flux between $h=300$ cm at 0.4 m and $h=100$ cm at 0 m, using Darcy's law $q=K(\Delta h/\Delta z-1)$ with the geometric mean of $K$.}\sol $K_r(100)=0.000324$, $K_r(300)=1.05e-05$, $K_r(1000)=2.26e-07$ \ans $K=0.00334$, $0.000108$, $2.33e-06$ cm\,d$^{-1}$ \ans Geometric mean of the two ends: $\bar K=\sqrt{0.00334\times0.000108}=0.0006$ cm\,d$^{-1}$ \ans Gradient: $\Delta h/\Delta z=(300-100)/40=5.0$, so $q=0.0006\times(5.0-1)=0.0024$ cm\,d$^{-1}$ \ans The flux is negligible against the demand of a crop (0.5 cm per day).

\exbox{Ten-soil example 2a: Clay loam (S09)}{Clay loam: $K_s=5.5$ cm\,d$^{-1}$ (typical for the class), $\alpha=0.0190$, $n=1.310$. Find $K(h)$ at 100, 300 and 1000 cm suction and the steady upward flux between $h=300$ cm at 0.4 m and $h=100$ cm at 0 m, using Darcy's law $q=K(\Delta h/\Delta z-1)$ with the geometric mean of $K$.}\sol $K_r(100)=0.00574$, $K_r(300)=0.000392$, $K_r(1000)=1.54e-05$ \ans $K=0.0317$, $0.00216$, $8.5e-05$ cm\,d$^{-1}$ \ans Geometric mean of the two ends: $\bar K=\sqrt{0.0317\times0.00216}=0.00828$ cm\,d$^{-1}$ \ans Gradient: $\Delta h/\Delta z=(300-100)/40=5.0$, so $q=0.00828\times(5.0-1)=0.0331$ cm\,d$^{-1}$ \ans The flux is negligible against the demand of a crop (0.5 cm per day).

\exbox{Ten-soil example 2a: Silty clay loam (S10)}{Silty clay loam: $K_s=3.6$ cm\,d$^{-1}$ (typical for the class), $\alpha=0.0100$, $n=1.230$. Find $K(h)$ at 100, 300 and 1000 cm suction and the steady upward flux between $h=300$ cm at 0.4 m and $h=100$ cm at 0 m, using Darcy's law $q=K(\Delta h/\Delta z-1)$ with the geometric mean of $K$.}\sol $K_r(100)=0.0139$, $K_r(300)=0.00153$, $K_r(1000)=8.64e-05$ \ans $K=0.0499$, $0.00551$, $0.000311$ cm\,d$^{-1}$ \ans Geometric mean of the two ends: $\bar K=\sqrt{0.0499\times0.00551}=0.0166$ cm\,d$^{-1}$ \ans Gradient: $\Delta h/\Delta z=(300-100)/40=5.0$, so $q=0.0166\times(5.0-1)=0.0663$ cm\,d$^{-1}$ \ans The flux is negligible against the demand of a crop (0.5 cm per day).

